\section{Preparação estrutural e bancos especializados}
A construção separa dados de origem, relações estruturais e candidatos prontos para consulta. O arquivo GTFS alimenta o grafo e as ligações diretas. A malha OSM, processada pelo OSRM com perfil pedestre, fornece distâncias verificadas entre paradas; essas distâncias complementam as ligações de ônibus. O fluxo de dependências é: cadastro GTFS, grafo, conexões diretas e caminhadas, componentes reutilizáveis e, por fim, materialização seletiva. Alterações das fontes ou dos parâmetros relevantes exigem verificar as assinaturas dos derivados e reconstruir os artefatos afetados, sem refazer essa estrutura estática a cada consulta operacional.

Os seis bancos SQLite especializados têm as seguintes responsabilidades. Os nomes abaixo indicam os sufixos usados pelo conjunto de artefatos \texttt{sptrans}; não incluem o banco Django do cenário histórico de Mauá.
\begin{enumerate}
\item \textbf{Cadastro GTFS: \nolinkurl{sptrans.sqlite3}.} Persiste paradas, linhas, viagens, sequências, horários, frequências, calendários, exceções e geometrias importadas. Chaves por identificador e índices das relações entre viagens, paradas e horários sustentam as consultas. É a fonte dos demais artefatos e da programação utilizada pelos dois planejadores.
\item \textbf{Grafo: \nolinkurl{sptrans.routing.sqlite3}.} Deriva do cadastro e organiza padrões de sequência, segmentos dirigidos, calendários e dados temporais para avaliação. As chaves de padrão e sequência permitem recuperar segmentos e ocorrências de paradas. Também comporta perfis históricos por segmento, tipo de dia e faixa horária e correspondências operacionais. A existência desse esquema não significa que uma base de velocidades históricas já tenha sido preenchida ou validada.
\item \textbf{Conexões diretas: \nolinkurl{sptrans.direct.sqlite3}.} Deriva do grafo e enumera, em cada padrão, embarques e desembarques posteriores na sequência, com distâncias e referências temporais. Os índices \nolinkurl{direct_out} e \nolinkurl{direct_in} permitem acesso pela origem e pelo destino. O índice reverso consulta ligações no sentido correto de circulação; não cria uma viagem no sentido contrário.
\item \textbf{Caminhadas: \nolinkurl{sptrans.walks.sqlite3}.} Depende das paradas do grafo e da malha pedestre OSM/OSRM. Persiste ligações verificadas, distâncias de rede e acessos à malha, tempo derivado da velocidade configurada, procedência e estado de verificação. As tabelas \nolinkurl{walk_edge} e \nolinkurl{stop_snap}, com índice de entrada \nolinkurl{walk_in}, apoiam a consulta nos dois sentidos. Distância e tempo permanecem distintos. A geometria completa pode ser recuperada posteriormente para apresentação, em vez de duplicada em cada candidato.
\item \textbf{Componentes: \nolinkurl{sptrans.connections.sqlite3}.} Combina conexões diretas e caminhadas em ocorrências, ligações pedestres, transferências e saídas de ônibus seguidas de caminhada. Índices de saída, entrada e encontro permitem reutilizar componentes nas composições de um, dois ou três ônibus. As visões \nolinkurl{conexoes_1}, \nolinkurl{conexoes_2} e \nolinkurl{conexoes_3} não significam armazenamento físico de todas as combinações completas.
\item \textbf{Candidatos materializados: \nolinkurl{sptrans.materialized.sqlite3}.} Deriva dos componentes e guarda, na tabela \nolinkurl{conexoes}, referências aos trechos e caminhadas, totais de distância, tempo a pé e classificação por categoria. Índices por par, distância e sentido reverso apoiam a recuperação. O estado de materialização é registrado por par origem--destino e número de embarques, distinguindo categorias concluídas de interrupções por limite de execução. Metadados vinculam os candidatos às fontes, parâmetros e cotas usados na preparação.
\end{enumerate}

A composição bidirecional ocorre offline: combina prefixos recuperados a partir da origem com sufixos consultados pelo destino, verificando encontros compatíveis nas sequências dirigidas. O módulo de três ônibus reutiliza esses componentes para evitar enumerar previamente todo produto cartesiano de paradas. A configuração de referência retém até 16 candidatos de um ônibus, oito de dois e quatro de três por par, segundo distância. Essa seleção reduz o conjunto avaliado online, mas pode excluir um percurso mais longo que seria mais rápido no horário consultado. A cobertura municipal e a perda de alternativas deverão ser medidas; um arquivo existente não comprova materialização concluída.
% TODO-BENCHMARK: Inserir tempo total da materialização.
% TODO-BENCHMARK: Inserir armazenamento final dos bancos.
% TODO-BENCHMARK: Inserir cobertura das paradas e pares.

\section{Consulta temporal e acionamento do fallback}
Com o seletor de planejador configurado para conexões e até duas baldeações, a consulta principal abre os artefatos preparados para leitura, verifica sua compatibilidade e identifica paradas candidatas de acesso e saída. A recuperação considera pares cujas três categorias de embarques estejam concluídas. Não exige que a cidade inteira esteja pronta para utilizar um par concluído. Os percursos são recuperados dos índices; a composição bidirecional completa não é repetida nessa etapa.

A classificação preliminar combina deslocamentos nos veículos, caminhadas e uma penalidade de referência de cinco minutos por espera intermediária. Essa penalidade serve à ordenação inicial, não substitui a espera calculada na avaliação temporal. As caminhadas de acesso, transferência e saída são validadas pelo serviço pedestre, respeitando o limite configurado. O fluxo materializado compartilha um orçamento de até 64 verificações e 30 segundos para iniciar novas verificações; os limites não asseguram que toda resposta termine em 30 segundos. Uma caminhada inválida impede publicar aquele candidato como percurso verificado.

O avaliador propaga o instante de chegada entre os trechos. Para cada embarque, consulta uma partida compatível com a chegada do passageiro, considerando caminhada e margem de transferência de 120 segundos. Calendários e exceções GTFS restringem os serviços; horários superiores a 24:00:00 preservam sua associação ao dia de serviço. A avaliação considera dias de serviço anteriores pertinentes e o seguinte, dentro do horizonte configurado de seis horas. O tempo total inclui acesso, espera inicial, permanência nos veículos, caminhadas e esperas nas transferências e saída até o destino. Distância curta, isoladamente, não assegura uma conexão temporal possível.

O mecanismo de custos admite perfis históricos de circulação por segmento, tipo de dia e faixa horária, com passagem entre faixas durante o percurso. Na ausência de um perfil completo aplicável ao segmento, utiliza o custo programado correspondente. Tempos de parada são tratados separadamente para evitar dupla contagem; velocidade comercial, que já os incorpora, não deve ser inserida como velocidade de circulação sem decomposição. Nesta revisão, não se apresenta resultado de uma base histórica preenchida. A infraestrutura de custos e a disponibilidade efetiva desses dados são estados distintos.

O refinamento com Olho Vivo exige correspondência de linha, sentido e parada e aceita previsões compatíveis com o instante de chegada ao embarque. A posição operacional não fornece automaticamente a velocidade de todos os segmentos nem um tempo exato da viagem inteira. Na indisponibilidade de previsão, permanece a avaliação programada; essa ocorrência isolada não aciona o planejador por rodadas. O orçamento operacional limita as consultas externas, sem reconstruir os artefatos estáticos. As alternativas válidas são reordenadas e enviadas progressivamente, retendo até cinco resultados distintos na configuração de referência. A duração continua sendo estimada e sua ordem pode mudar durante o processamento.

O despachante aciona a busca por rodadas quando a preparação não retorna contexto utilizável, inclusive ausência de candidatos ou incompatibilidade detectada dos artefatos, ou quando o resultado final materializado não contém rotas válidas. As falhas de artefato tratadas pelo adaptador também podem conduzir a esse retorno vazio. Se o fluxo materializado termina com alternativas, a consulta encerra por esse caminho, ainda que alguma limitação tenha ocorrido: não executa automaticamente o fallback para completar a lista. A seleção explícita do planejador por rodadas, ou uma profundidade fora do limite suportado pelo caminho materializado, conduz diretamente ao mecanismo alternativo. Não se afirma recuperação automática de qualquer erro de programação.

Os dois planejadores não são executados simultaneamente nem têm seus resultados obrigatoriamente unidos. O identificador do produtor acompanha o retorno; a mensagem atual de fallback é genérica. Para os benchmarks, será necessário registrar motivos discriminados no instrumento de medição, sem deduzi-los apenas dessa mensagem. As cotas por distância, os limites de validação e o filtro de alternativas distintas não constituem uma fronteira de Pareto completa. As regras de poda por chegada descritas a seguir pertencem ao planejador por rodadas; não se atribui à consulta materializada uma poda global pelo melhor tempo que ela não executa.

\section{Protocolo do benchmark definitivo}\label{sec:benchmark-hibrido}
O protocolo será executado após a conclusão e verificação da materialização. Serão comparados: A, planejador materializado isolado, com resultado vazio mantido como tal; B, planejador por rodadas isolado; e C, arquitetura integrada com fallback. Cada resultado deverá identificar o mecanismo que efetivamente o produziu. O plano não constitui resultado experimental e não exige executar benchmarks durante esta revisão do relatório.

As três configurações utilizarão os mesmos pares origem--destino, instantes de partida e fuso, pacote e calendários GTFS, velocidade de caminhada de 4,5 km/h, limites de distância, número máximo de embarques, critérios temporais e informações operacionais. O instrumento deverá equalizar também o tratamento de caminhadas iniciais e finais e o refinamento temporal; apenas aplicar os parâmetros aos fluxos existentes não comprova equivalência. Serão registrados a revisão do código, eventuais alterações locais, assinaturas dos bancos, configuração, hardware e versões dos serviços. A amostra deverá incluir viagens diretas, baldeações, caminhadas de acesso e saída, períodos diurnos e noturnos e casos sem rota. Selecionar somente pares com sucesso no materializado impediria medir sua cobertura e a necessidade do fallback.

\textbf{Teste controlado.} As respostas operacionais e pedestres serão congeladas e reutilizadas. As consultas serão repetidas em execuções isoladas, alternando a ordem dos planejadores. Serão separadas condições de cache frio e quente, documentando quais caches foram reiniciados ou preservados, inclusive os do sistema operacional. O número de repetições e a amostra serão definidos antes da coleta; percentis extremos exigem quantidade suficiente de observações. Serão medidos latência até a primeira alternativa e até o término, memória e atividade de preparação separadamente. A materialização e outros trabalhos intensivos não deverão competir com as consultas medidas.

\textbf{Teste operacional.} Os planejadores serão consultados sequencialmente em instantes próximos, com dados atuais da SPTrans, alternando sua ordem e registrando os instantes de requisição e das informações retornadas. A comparação examinará viabilidade dos embarques, alternativas e sensibilidade à atualização dos dados. Respostas diferentes entre duas consultas serão explicitadas, não tratadas como se fossem a mesma observação congelada. Eventual concorrência por recursos ou variação da rede não será interpretada como diferença intrínseca entre algoritmos.

% TODO-BENCHMARK: Inserir latências até a primeira alternativa e até o término, por configuração e cache.
% TODO-BENCHMARK: Inserir qualidade dos itinerários, viabilidade e diferenças temporais.
% TODO-BENCHMARK: Inserir frequência, motivos e custo do fallback no fluxo integrado.
% TODO-BENCHMARK: Inserir comparação A (materializado), B (rodadas) e C (integrado), com amostra e limites.
\textbf{Métricas e critérios de análise.} Para cada configuração e condição de cache, registrar latências p50, p95 e p99, média, desvio-padrão e tamanho da amostra; quantidade de candidatos examinados e alternativas temporalmente válidas; melhor duração, baldeações, caminhada e espera, incluindo sua decomposição por etapa. Registrar consumo de memória, tamanho final de cada banco e do conjunto e tempo total de preparação, discriminando suas etapas. No fluxo integrado, medir frequência e motivos de fallback, tempo consumido antes de seu acionamento e produtor de cada resposta.

A qualidade será comparada por cobertura, compatibilidade temporal e diferença entre as melhores durações encontradas, sem assumir que um dos planejadores limitados fornece o ótimo global. Deverá ser examinada a perda causada pelas cotas de 16, oito e quatro candidatos, mediante amostras com cotas ampliadas ou referência independente, quando viável. Resultados sem rota serão contabilizados e terão suas causas distinguidas de timeout e falha externa. A conclusão sobre o compromisso entre preparação, armazenamento, latência e alternativas somente será redigida após essas medições.


\textbf{Comparabilidade das alternativas.} Para não confundir igualdade de representação com utilidade equivalente, o protocolo distinguirá: itinerários com a mesma sequência dirigida de linhas, paradas e ocorrências de embarque; o mesmo deslocamento com variações nos acessos a pé; itinerários distintos com chegada temporalmente equivalente; alternativa mais rápida encontrada por apenas um mecanismo; alternativa inviável; e ausência de resultado. A tolerância temporal e o critério para agrupar variações de caminhada deverão ser definidos antes da coleta, com preservação dos valores exatos de chegada, distância, espera e baldeações. Uma alternativa inviável não será classificada como equivalente por apresentar duração semelhante. Essa classificação complementa as métricas de cobertura e latência, sem alterar as três configurações comparadas ou antecipar resultados.

\section{Situação dos benchmarks em 23 de outubro de 2026}
A revisão dos artefatos disponíveis não localizou uma nova execução do benchmark definitivo. Os resultados encontrados nos diretórios de evidências datam de 7 e 8 de outubro e não foram promovidos a resultados finais. Foi realizado um inventário somente de leitura dos seis bancos utilizados pela aplicação, com registro de tamanho e SHA-256. O script e os dados brutos estão em \nolinkurl{benchmarks/2026-10-23/}. Trata-se de diagnóstico de preparação, não de comparação experimental dos planejadores.

\begin{table}[H]\centering
\caption{Armazenamento observado em 23 de outubro, antes da conclusão da materialização.}\label{tab:inventario-benchmark}
\begin{tabular}{lr}\hline
Arquivo & Tamanho (MiB)\\\hline
\nolinkurl{sptrans.sqlite3} & 70,04\\
\nolinkurl{sptrans.routing.sqlite3} & 595,36\\
\nolinkurl{sptrans.direct.sqlite3} & 402,20\\
\nolinkurl{sptrans.walks.sqlite3} & 446,32\\
\nolinkurl{sptrans.connections.sqlite3} & 2907,56\\
\nolinkurl{sptrans.materialized.sqlite3} & 112,23\\
\hline
Total & 4533,70\\\hline
\end{tabular}
\fonte{Inventário local somente de leitura; 1 MiB = 1.048.576 bytes.}
\end{table}

Os valores representam o tamanho lógico dos arquivos principais naquele instante, sem incluir eventuais arquivos auxiliares, caches ou exportações; não representam armazenamento final da materialização municipal. As assinaturas permaneceram estáveis quanto ao tamanho e à data de modificação durante a leitura de cada arquivo, sem constituir um snapshot transacional conjunto.

O banco materializado informa \texttt{city\_complete=false}. Seu catálogo contém 77.957 pares na categoria de um ônibus e 77.956 na categoria de dois ônibus, todos marcados como \texttt{ready}; 77.956 pares possuem ambas as categorias. Há 644.702 candidatos armazenados. O esquema desse artefato admite apenas essas duas categorias, não a categoria de três ônibus prevista no protocolo. O estado \texttt{ready} indica processamento da categoria, inclusive quando ela não produz caminho: as contagens não equivalem a cobertura de viagens viáveis. Os validadores existentes aceitaram as assinaturas dos componentes e do materializado, o que não comprova conclusão municipal nem equivalência das três configurações experimentais.

Permanecem sem medição definitiva o tempo total de preparação, a latência até a primeira alternativa e até o término, a cobertura temporal, a qualidade dos itinerários e a frequência e o custo de fallback. Não foram calculados percentis com base nos registros históricos nem somados tempos parciais como se representassem o tempo total de preparação. A comparação A/B/C deverá ocorrer após conclusão e verificação da materialização, com amostra, repetições, tolerâncias e condições de cache definidas previamente, conforme o protocolo preservado. Nenhuma materialização ou consulta comparativa foi iniciada nesta revisão.

